% Guion oral - Steven Pozo - 3 minutos
% Avance Día 1 - 20 de noviembre de 2026

\section*{Guion: arquitectura y modelo de comunicación}

\textbf{0:00--0:20. Alcance.}
Este diagrama no pretende mostrar toda la infraestructura de Netflix. Es una
vista lógica de un caso que Netflix sí documentó: su plataforma de comprensión
de medios. Así evitamos inventar nombres de servicios o conexiones internas.

\textbf{0:20--1:05. Componentes.}
A la izquierda separo dos consumidores. La aplicación o interfaz consumidora
entra por la ruta GraphQL documentada. El servicio backend consumidor es el
cliente gRPC: posee el stub y realiza la RPC unaria. Esta separación evita
afirmar, sin evidencia, que la aplicación final llama directamente a gRPC.
En el límite de la plataforma existen dos interfaces documentadas: GraphQL para
integraciones de interfaz y gRPC, usado principalmente de backend a backend.
La ruta que nos interesa llega al Search Gateway. Este módulo coordina el
procesamiento de la consulta y de la respuesta y se conecta con capacidades de
búsqueda e índices de datos de medios. Las cajas representan responsabilidades
lógicas, no máquinas específicas ni cantidades de réplicas.

\textbf{1:05--1:45. Cliente-servidor y RPC unaria.}
La aplicación es el cliente y el servicio gRPC es el servidor. Ambos comparten
un contrato definido con Protocol Buffers. El cliente llama un método en un
stub; gRPC serializa y transporta el mensaje por HTTP/2. En una RPC unaria el
cliente envía exactamente una petición y recibe una respuesta con datos o un
estado de error. Se parece a una función local, pero sigue siendo una llamada de
red y por eso puede tener latencia, vencimiento o fallo.

La secuencia se puede recordar con cuatro verbos: \textbf{invocar, serializar,
ejecutar y responder}. Primero se invoca el stub; luego se serializa la petición;
el servidor ejecuta la operación; finalmente devuelve la respuesta y el estado.
Si no llega a tiempo, el cliente no debe asumir que fue una llamada local exitosa.

\textbf{1:45--2:25. Por qué gRPC.}
La evidencia de Netflix indica dos necesidades: comunicación entre servicios
escritos en diferentes lenguajes y evitar bibliotecas cliente mantenidas a
mano. gRPC genera código desde el contrato y da una interfaz uniforme. En el
caso de la plataforma de medios, Netflix afirma explícitamente que gRPC se usa
principalmente entre backends; esa es la evidencia que sostiene la flecha del
diagrama.

\textbf{2:25--3:00. Límite y cierre.}
gRPC resuelve la invocación remota, pero no garantiza por sí solo que Netflix
sea disponible, consistente o tolerante a fallos. El diagrama no atribuye
replicación, balanceo, reintentos ni descubrimiento sin fuente. La idea clave es
separar el contrato de comunicación de la topología física y de los mecanismos
operativos.

\subsection*{Preguntas probables}

\begin{itemize}
  \item \textbf{¿Netflix usa gRPC en todo?} No. La fuente dice que predomina en
  backend a backend; la misma plataforma también expone GraphQL.
  \item \textbf{¿Una RPC es local?} No. El stub oculta parte de la mecánica,
  pero la llamada cruza la red y puede fallar.
  \item \textbf{¿Dónde entra RMI?} Como alternativa Java--Java para comparar;
  no hay evidencia presentada aquí de que Netflix use Java RMI.
  \item \textbf{¿Protocol Buffers y gRPC son lo mismo?} No. Protocol Buffers
  define y serializa mensajes; gRPC organiza la llamada remota y sus estados.
  \item \textbf{¿Por qué el diagrama no muestra regiones o réplicas?} Porque las
  fuentes utilizadas para este caso no demuestran esa topología física concreta.
\end{itemize}
